%
% reversi_070_feature.tex
%
% This file is part of the reversi program
% http://github.com/rcrr/reversi-v2
%
% Author Roberto Corradini mailto:rob_corradini@yahoo.it
% Copyright 2026 Roberto Corradini. All rights reserved.
%
% License
%
% This program is free software; you can redistribute it and/or modify it
% under the terms of the GNU General Public License as published by the
% Free Software Foundation; either version 3, or (at your option) any
% later version.
%
% This program is distributed in the hope that it will be useful,
% but WITHOUT ANY WARRANTY; without even the implied warranty of
% MERCHANTABILITY or FITNESS FOR A PARTICULAR PURPOSE. See the
% GNU General Public License for more details.
%
% You should have received a copy of the GNU General Public License
% along with this program; if not, write to the Free Software
% Foundation, Inc., 59 Temple Place - Suite 330, Boston, MA  02111-1307, USA
% or visit the site <http://www.gnu.org/licenses/>.
%
\chapter{Le feature}
\label{ch:feature}

Fino a questo capitolo abbiamo descritto una posizione di gioco con due strumenti: la
board, i 64 quadrati e la loro occupazione, e i pattern, le configurazioni locali
canonizzate per simmetria. Il modello che costruiremo però non si alimenta di pattern
da soli: ha bisogno di un concetto più generale, quello di \textit{feature}.

Una \textbf{feature} è una caratteristica osservabile di una posizione di gioco: è la
variabile che, letta sulla board, viene offerta in input al modello e ne guida
l'inferenza. I pattern sono la fonte di feature più ricca e sono stati trattati nel
capitolo che li precede; ma non sono l'unica. A essi si affiancano due feature di
natura diversa, che completano il quadro: la \textit{mobility}, una misura delle
opzioni di movimento dei due giocatori, e l'\textit{intercetto}, il termine costante
del predittore. Il \textit{FeatureSet} di una partita è la raccolta nominata di
feature scelte a piacere tra queste tre fonti, ed è ciò che determina, nella fase di
setup, le colonne della matrice di disegno che il modello impara a pesare.

\section{Le tre categorie di feature}

Ogni feature avvolge una delle tre \textit{fonti di feature}, individuate da una
categoria:

\begin{itemize}
    \item \textbf{PATTERN} --- una selezione di $S$ quadrati sulla board, canonizzata
          per simmetria $D_4$ e indirizzata dal proprio rappresentante principale. È la
          fonte più espressiva: un pattern può presentare molte configurazioni (le
          assegnazioni base-3 dei suoi quadrati) e apparire in più istanze sulla board.
    \item \textbf{MOBILITY} --- una misura scalare delle mosse disponibili per un
          giocatore, prese su un insieme di quadrati (la maschera) e lette dal punto di
          vista di uno dei due colori.
    \item \textbf{INTERCEPT} --- il termine costante: una feature che non osserva
          nulla sulla board e vale uguale in ogni posizione.
\end{itemize}

La differenza strutturale tra le tre fonti è il modo in cui una feature si scompone in
\textit{istanze} e \textit{configurazioni}. Un pattern ha molte istanze (una per ogni
occorrenza sulla board) e molte configurazioni (una per ogni assegnazione base-3). Una
mobility ha un'unica istanza ma molte configurazioni (una per ogni possibile valore del
conteggio). L'intercetto ha un'unica istanza e un'unica configurazione. È questa
decomposizione in istanze $\times$ configurazioni a definire il numero di colonne che
ogni feature contribuisce al modello.

Un \textbf{FeatureSet} è una raccolta nominata e ordinata di feature, scelte a piacere
tra intercetto, mobility e pattern. L'ordine è significativo: le colonne della matrice
di disegno sono disposte nel gruppo per feature, e all'interno di ogni feature nel
gruppo per istanza e per configurazione. Cambiare il FeatureSet cambia la forma del
modello senza cambiarne la natura: è lo strumento con cui si esplora il compromesso tra
ricchezza descrittiva e numero di parametri da stimare.

\section{L'intercetto}

L'intercetto è la feature costante del predittore: il termine $b$ che accompagna la
somma dei contributi di pattern e mobility. È una feature di una sola istanza e una
sola configurazione --- non osserva alcunché sulla board, ma fornisce al modello il
\textit{bias}, il livello di partenza da cui i contributi osservati prendono le mosse.
Nel software si crea con \texttt{Feature.new\_intercept()} (categoria
\texttt{INTERCEPT}); nel FeatureSet è opzionale, ed è controllato dalla flag
\texttt{intercept} della configurazione. In termini di matrice di disegno,
l'intercetto è la colonna di tutti uni: contribuisce con un unico peso, uguale per
ogni riga (ogni posizione osservata), e assorbe la componente media della variabile
target che le altre feature, centrate, non sono in grado di rappresentare.

\section{La mobility}

La mobility è la feature che misura le opzioni di movimento: quante mosse legali un
giocatore ha a disposizione, sull'intera board o su una sua sotto-regione. È il
segnale strategico più immediato di un gioco a rincorsa come il Reversi. In una
posizione dimostra, in un colpo d'occhio, chi dei due colori è in grado di agire e chi
è costretto a reagire: un giocatore le cui mosse si assottigliano è un giocatore che
sta perdendo il controllo della board, e lo squilibrio tra le mobility dei due colori
è uno dei predittori più forti --- e più parsimoniosi in parametri --- dell'esito
della partita.

A differenza di un pattern, che è una \textit{forma} geometrica cercata sulla board,
la mobility è un \textit{numero}: il conteggio di quante mosse sono possibili. È un
osservabile globale della posizione, definito da due soli ingredienti --- una
\textit{maschera}, che stabilisce \textit{dove} contare, e una \textit{bandiera di
lato}, che stabilisce \textit{dal punto di vista di chi} contare. È da questi due
ingredienti che derivano sia la forma formale della feature sia il suo significato
strategico, e li prendiamo in ordine.

\subsection{Definizione formale}

Formalmente, la mobility è una feature a \textit{table-lookup} scalare. Data una
posizione, il motore di gioco calcola due insiemi di mosse: le mosse legali del
giocatore di turno, $lms$ (\textit{legal moves}), e le mosse che l'avversario avrebbe
se fosse il suo turno, $alms$ (\textit{anti-legal moves}, le cosiddette
\textit{would-be moves}). La maschera $\mathcal{M}$ è un \texttt{Bitboard} che
seleziona i quadrati su cui il conteggio viene preso; l'indice della feature è il
numero di quadrati nell'intersezione dell'insieme di mosse con la maschera:

\begin{equation}
\text{indice} \;=\;
\begin{cases}
\operatorname{popcount}\!\left(lms \,\&\, \mathcal{M}\right) & \text{se } \texttt{side} = \text{mover} \\[4pt]
\operatorname{popcount}\!\left(alms \,\&\, \mathcal{M}\right) & \text{se } \texttt{side} = \text{opponent}
\end{cases}
\end{equation}

La maschera gioca qui un ruolo affine a quello che ha per i pattern, ma con una
differenza essenziale: per la mobility non viene canonizzata. La maschera è un
selettore di quadrati --- l'insieme dei quadrati su cui si vuole misurare
l'attività --- e due maschere che sono l'una l'immagine dell'altra per una simmetria
$D_4$ restano due feature distinte, perché la mobility non ha un'idea di
``rappresentante principale''. È la maschera a dare alla mobility la sua
\textit{granularità}: con la maschera piena, tutti e 64 i quadrati, la mobility conta
le mosse sull'intera board ed è un indicatore \textit{globale}; con una maschera
ridotta --- solo i quadrati centrali, solo i bordi, solo i quadrati adiacenti a un
angolo --- la mobility diventa un indicatore \textit{localizzato} dell'attività in una
regione. Una singola posizione può così alimentare non una ma diverse mobility, una per
ogni maschera di interesse: la famiglia delle mobility è indicizzata sulle maschere,
esattamente come la famiglia dei pattern è indicizzata sulle forme.

La bandiera \texttt{side} sceglie il punto di vista dal quale il conteggio è letto. Le
regole del Reversi definiscono le mosse disponibili solo per il giocatore di turno: per
\texttt{side} = \texttt{mover} la mobility conta le mosse che il giocatore attuale può
giocare. Il \texttt{side} = \texttt{opponent} è invece una definizione a livello di
feature, non di gioco: conta le mosse che l'avversario \textit{avrebbe} se fosse il suo
turno. È proprio questa lettura ``dall'altro lato'' a rendere la mobility una feature
così informativa: non misura solo quante mosse ho io, ma quante ne sto concedendo
all'avversario.

\subsection{Le due mobility standard: LMC e ALMC}

I due casi standard usano la maschera piena e differiscono solo per la bandiera di
lato:

\begin{itemize}
    \item \textbf{LMC} --- \textit{Legal Move Count}: $\operatorname{popcount}(lms)$,
          il numero di mosse legali del giocatore di turno;
    \item \textbf{ALMC} --- \textit{Anti Legal Move Count}:
          $\operatorname{popcount}(alms)$, il numero di mosse che l'avversario avrebbe
          al suo turno.
\end{itemize}

La coppia LMC/ALMC è la mobility di riferimento del modello, e va letta come una
coppia, non come due feature scollegate. Presa da sola, la LMC dice quanto è
``aperta'' la posizione per chi deve muovere; presa con la ALMC, dice anche quanto la
mossa scelta sta restringendo --- o allargando --- lo spazio di manovra
dell'avversario. C'è qui una lettura che va oltre il semplice conteggio: la ALMC è la
misura delle mosse che la mia mossa \textit{lascia} all'avversario, e un valore alto
di ALMC è il segno di una mossa ``generosa'', che allarga il gioco. Molte delle mosse
che un giocatore esperto giudica ``forti'' nel centro partita sono mosse che
massimizzano la propria LMC mentre minimizzano la ALMC: la mobility fa esplicito, in
due numeri, quello che l'istinto del giocatore esprime in termini di ``controllo'' e
``respiro''.

\subsection{Un esempio a mano}

Fissiamo una posizione di centro partita e calcoliamo a mano i valori, per rendere
concreti i simboli. Supponiamo che il giocatore di turno, nero, abbia 12 mosse legali
e che l'avversario, bianco, ne avrebbe 9 se fosse il suo turno. Con la maschera piena,
l'indice della LMC è $\operatorname{popcount}(lms) = 12$ e quello della ALMC è
$\operatorname{popcount}(alms) = 9$: il FeatureSet osserva, per questa posizione, la
coppia di valori $(12, 9)$, che andranno a indizzare la rispettiva tabella di pesi.
Se lo stesso stato di board fosse ruotato di 90 gradi, i due conteggi resterebbero
identici --- 12 e 9 --- perché il \texttt{popcount} non cambia per rotazione: è questa
invarianza, a cui torneremo parlando di istanze, a fare della mobility una feature a
istanza singola.

Ora immaginiamo di restringere la maschera ai sedici quadrati centrali della board,
quelli che nel gergo del gioco formano il ``centro''. La stessa posizione produce
allora una \textit{mobility di centro}: l'indice è $\operatorname{popcount}(lms \,\&\,
\mathcal{M}_{\text{centro}})$, e vale, diciamo, 5 per il nero e 4 per il bianco. I due
numeri $(5, 4)$ dicono quante delle mosse disponibili cadono nel centro --- un
indicatore localizzato che la LMC/ALMC globale, aggregando su tutti e 64 i quadrati,
non può dare. È così che, a partire dalla stessa definizione, si ottiene una famiglia
di mobility con risoluzioni diverse: globale, di centro, di bordo, di angolo.

\subsection{Configurazioni, istanze e tabella dei pesi}

Come per i pattern, la mobility è una feature a \textit{table-lookup}: l'indice
calcolato sopra indizza una tabella di pesi. Se la maschera copre $k =
\operatorname{popcount}(\mathcal{M})$ quadrati, l'indice varia in $0, 1, \ldots, k$ e
ogni valore di questo intervallo è raggiungibile da qualche posizione, dunque il numero
di configurazioni della feature è $n_{\text{config}} = k + 1$, che è anche la
dimensione della sua tabella di pesi. Con la maschera piena ($k = 64$) la tabella ha 65
entrate, una per ogni possibile conteggio di mosse da 0 a 64; con la maschera di centro
($k = 16$) ne ha 17. Ogni entrata è un parametro del modello: il training impara, per
ogni possibile valore della mobility, quanto quel valore contribuisce alla valutazione
della posizione.

Qui sta la differenza strutturale più importante rispetto ai pattern, ed è il contrario
di quella che si potrebbe intuire. Un pattern ha \textit{molte istanze} e
\textit{molte configurazioni}; una mobility ha \textit{una sola istanza} e
\textit{molte configurazioni}. La mobility non ha istanze multiple:
$n_{\text{istanze}} = 1$. La coppia (maschera, lato) definisce la feature
direttamente, e non esiste un'orbita di simmetrie da canonizzare. La ragione è che il
conteggio di \texttt{popcount} è di per sé invariante per simmetria: ruotare o
riflettere la board non cambia \textit{quante} mosse ci sono, cambia solo
\textit{dove} sono. Poiché la mobility non guarda al ``dove'' ma al ``quante'', la
simmetria non le crea istanze distinte. Una stessa maschera letta con i due valori di
\texttt{side} è invece una coppia di feature distinte --- la LMC e la ALMC sono due
feature, non due istanze di una stessa feature --- e questa è l'unica
``moltiplicazione'' che la mobility conosce.

A differenza delle configurazioni di un pattern --- che codificano \textit{chi} sta su
ogni quadrato del pattern, con la loro combinatoria base-3 --- le configurazioni della
mobility codificano \textit{quante} mosse ci sono: è una feature ordinale, il cui
ordine naturale è quello numerico del conteggio. Il modello non impara un peso
separato per ogni ``forma'' di mobility (le forme, per la mobility, non esistono), ma
un peso per ogni possibile \textit{valore} del conteggio: una scala di pesi, dal caso
``nessuna mossa'' al caso ``mossa su ogni quadrato della maschera''.

\subsection{La mobility in confronto alle altre feature}

La tabella riassume, in una vista d'insieme, come le tre fonti di feature si
distribuiscono sulla griglia istanze $\times$ configurazioni, e cosa ciascuna osserva
sulla board:

\begin{table}[ht]
\centering
\begin{tabular}{lccc}
\toprule
 & \textbf{PATTERN} & \textbf{MOBILITY} & \textbf{INTERCEPT} \\
\midrule
osserva        & una forma locale & un conteggio & nulla \\
istanze        & molte & una & una \\
configurazioni & molte & molte & una \\
canonizzazione & sì & no & --- \\
\bottomrule
\end{tabular}
\caption{Le tre fonti di feature sulla griglia istanze $\times$ configurazioni.}
\label{tab:feature-sources}
\end{table}

Leggendo la tabella per righe si misura il ``peso descrittivo'' di ciascuna fonte: il
pattern è la feature più ricca (molte istanze e molte configurazioni, compattate dalla
canonizzazione), la mobility è una feature intermedia (una sola istanza, ma una scala
ampia e ordinata di configurazioni), e l'intercetto è la feature minima (una riga e
una colonna). È la \textit{combinazione} delle tre, nel FeatureSet, a dare al modello
la sua piena capacità descrittiva.

\subsection{La mobility come segnale strategico}

Perché una feature così semplice --- un conteggio di mosse --- merita uno spazio
proprio? Perché in un gioco a rincorsa la mobilità \textit{è} la dinamica del gioco.
Nella fase centrale le mobility dei due colori oscillano: una mossa che ``stringe''
l'avversario (una chiusura, un \textit{block} su una zona della board) abbassa la ALMC
e inasprisce, di rimando, la pressione sulle mosse successive. Il modello, osservando
la sequenza di valori di LMC e ALMC lungo la partita, impara a leggere queste
oscillazioni come indicatori di chi sta prendendo il controllo.

In prossimità della fine partita la mobility diventa il segnale dominante. Quando le
mosse di un giocatore su una regione scendono a zero, quella regione è di fatto
chiusa: non ci sarà più gioco lì, e il suo destino è già scritto. Lo squilibrio di
mobility anticipa così il \textit{freeze} finale --- quel momento in cui la board si
``gela'' e le mosse rimanenti sono forzate --- con diverse mosse di anticipo. Per
questo il modello tratta la mobility come una feature ordinale a tutti gli effetti: il
valore del conteggio è l'osservabile, e il peso che il training associa a ogni valore
è la stima imparata di quanto quel valore conti, in termini di esito, in quel punto
della partita.

Insieme, i pattern (le feature locali, ``chi è dove''), la mobility (la feature
globale ordinale, ``quante mosse'') e l'intercetto (il bias costante) formano il
FeatureSet che il predittore additivo del capitolo successivo impara a pesare. La
mobility ci porta così dalla descrizione geometrica della board alla sua lettura in
termini di gioco: è il ponte tra la teoria dei pattern e la funzione di valutazione.
